\documentclass[conference]{IEEEtran}

\usepackage[utf8]{inputenc}
\usepackage[T1]{fontenc}
\usepackage{newtxtext,newtxmath}
\usepackage{amsmath,amssymb,bm}
\usepackage{mathtools}
\usepackage{booktabs}
\usepackage{siunitx}
\usepackage{enumitem}
\usepackage[breaklinks=true,hidelinks]{hyperref}

\title{A Practical Description of the PLL Frequency Tracker\\
and Its C\texttt{++} Implementation}

\author{%
  \IEEEauthorblockN{Mikhail S.\ Grushinskiy}
  \IEEEauthorblockA{Independent Researcher\\
  Bareboat Necessities Project\\
  \texttt{mgrouch@users.github.com}}
}

\begin{document}
\maketitle

\begin{abstract}
This note gives a clean and practical description of the frequency tracker
implemented in \texttt{PLLFreqTracker}. The method is a lightweight
PLL/FLL-style tracker for one dominant narrowband component. It combines
front-end drift removal, synchronous demodulation, low-pass baseband
filtering, and a bounded PI loop on oscillator frequency.

The implementation is designed for variable sample intervals, explicit
frequency-band limits, confidence and lock logic, and an optional coarse
acquisition aid. The main intended input is world-frame vertical inertial
acceleration with gravity removed, although any consistent
acceleration-like units may be used.
\end{abstract}

\begin{IEEEkeywords}
frequency estimation, phase-locked loop, synchronous demodulation, wave
tracking, embedded signal processing
\end{IEEEkeywords}

% ---------------------------------------------------------------
\section{Introduction}

The tracker is intended to follow one dominant oscillatory component in a
signal such as vertical inertial acceleration. Its internal structure is
simple:

\begin{enumerate}[leftmargin=1.5em]
  \item remove very slow drift,
  \item suppress high-frequency junk,
  \item demodulate against an internal oscillator,
  \item low-pass the demodulated baseband terms,
  \item form a phase error from \(\operatorname{atan2}(Q,I)\),
  \item update oscillator frequency with a bounded PI loop.
\end{enumerate}

Compared with a Kalman-style tracker, this implementation is explicitly
control-oriented. It does not carry a covariance matrix. Instead, it uses a
phase detector, loop gains, signal-quality logic, and several safety limits.

The design assumes one dominant narrowband component. If the signal is
strongly multimodal, any single-frequency tracker may hop between modes or
settle somewhere between them.

% ---------------------------------------------------------------
\section{Configuration and Main Outputs}

The tracker is configured by a structure \texttt{Config} containing:
\begin{itemize}[leftmargin=1.4em]
  \item a physical search band \((f_{\min}, f_{\max})\),
  \item an initial nominal frequency \(f_{\mathrm{init}}\),
  \item front-end high-pass and low-pass cutoffs,
  \item demodulation baseband low-pass cutoff,
  \item loop bandwidth and damping,
  \item smoothing time constants for output, power, and confidence,
  \item lock thresholds,
  \item slew-rate, recentering, and internal substep limits,
  \item optional coarse-acquisition parameters.
\end{itemize}

The main public outputs are:
\begin{itemize}[leftmargin=1.4em]
  \item \texttt{getFrequencyHz()} --- smoothed tracked frequency,
  \item \texttt{getRawFrequencyHz()} --- raw loop frequency before output smoothing,
  \item \texttt{getPhaseRad()} --- internal oscillator phase,
  \item \texttt{getPhaseErrorRad()} --- current phase-detector error,
  \item \texttt{getBandpassedSignal()} --- prefiltered signal used by the loop,
  \item \texttt{getAmplitudeEstimate()}, \texttt{getRmsEstimate()},
        \texttt{getCoherence()}, \texttt{getConfidence()},
  \item \texttt{isLocked()},
  \item \texttt{getCoarseFrequencyHz()} when coarse assist is valid.
\end{itemize}

\begin{table}[t]
  \centering
  \caption{Main default configuration values}
  \label{tab:defaults}
  \small
  \begin{tabular}{@{}lll@{}}
    \toprule
    Parameter & Default & Meaning \\
    \midrule
    \(f_{\min}\) & \SI{0.045}{Hz} & lower search band edge \\
    \(f_{\max}\) & \SI{0.35}{Hz} & upper search band edge \\
    \(f_{\mathrm{init}}\) & \SI{0.12}{Hz} & reset / nominal frequency \\
    \(f_{\mathrm{hp}}\) & \SI{0.015}{Hz} & drift-removal HP cutoff \\
    \(f_{\mathrm{lp}}\) & \SI{0.45}{Hz} & front-end LP cutoff \\
    \(f_{\mathrm{dem}}\) & \SI{0.05}{Hz} & baseband LP cutoff \\
    \(f_{\mathrm{loop}}\) & \SI{0.018}{Hz} & loop bandwidth \\
    \(\zeta\) & 1.0 & loop damping \\
    \(\tau_{\mathrm{out}}\) & \SI{4}{s} & output smoothing \\
    \(\tau_{\mathrm{pow}}\) & \SI{14}{s} & power smoothing \\
    \(\tau_{\mathrm{conf}}\) & \SI{10}{s} & confidence smoothing \\
    \(r_{\min}\) & 0.012 & lock RMS floor in input units \\
    \(c_{\mathrm{lock}}\) & 0.45 & lock coherence threshold \\
    \(c_{\mathrm{unlock}}\) & 0.25 & unlock coherence threshold \\
    \bottomrule
  \end{tabular}
\end{table}

% ---------------------------------------------------------------
\section{Reset State and Frequency Variables}

On reset, the tracker clamps the requested initial frequency into the valid
band and stores:
\begin{align}
  \omega_0 &= 2\pi f_{\mathrm{init}},\\
  \omega_{\mathrm{nom}} &= 2\pi f_{\mathrm{cfg}},
\end{align}
where \(f_{\mathrm{cfg}}\) is the configured nominal reset frequency.

The internal oscillator state is then initialized as
\begin{equation}
  \omega = \omega_0,
  \qquad
  \phi = 0.
\end{equation}

The public outputs are initialized as
\begin{equation}
  f_{\mathrm{raw}} = f_0,
  \qquad
  f_{\mathrm{smooth}} = f_0,
\end{equation}
where \(f_0\) is the reset frequency in hertz.

All filter states, confidence states, lock flags, and coarse-assist states are
also reset.

% ---------------------------------------------------------------
\section{Variable-\(dt\) Update and Internal Substeps}

The public update method receives one new input sample \(x\) and a sample
interval \(\Delta t\). If either value is non-finite, or if \(\Delta t \le 0\),
the update is skipped.

To improve numerical robustness under irregular timing, large updates are split
into smaller internal steps. If the configured maximum internal step is
\(\Delta t_{\max}\), then the code uses
\begin{equation}
  N = \left\lceil \frac{\Delta t}{\Delta t_{\max}} \right\rceil,
  \qquad
  \delta t = \frac{\Delta t}{N},
\end{equation}
and applies the internal update \(N\) times using the same input sample.

This is a distinctive implementation feature: the loop is written to behave
well when \(\Delta t\) varies, rather than assuming a fixed-rate sample clock.

% ---------------------------------------------------------------
\section{Front-End Preprocessing}\label{sec:frontend}

\subsection{Exact One-Pole Smoothing}

All one-pole filters in the implementation use the exact exponential form
\begin{equation}
  z \leftarrow \alpha z + (1-\alpha)u,
\end{equation}
rather than a small-\(\Delta t\) Euler approximation.

For a cutoff frequency \(f_c\), the coefficient is
\begin{equation}
  \alpha_f = \exp(-2\pi f_c \Delta t),
\end{equation}
and for a time constant \(\tau\), the coefficient is
\begin{equation}
  \alpha_\tau = \exp\!\left(-\frac{\Delta t}{\tau}\right).
\end{equation}

This choice keeps the one-pole stages stable for any \(\Delta t > 0\).

\subsection{Drift Removal and Broad Bandpass}

The first preprocessing stage removes very slow drift by forming a high-pass
signal as input minus a low-pass estimate:
\begin{align}
  d &\leftarrow \alpha_{\mathrm{hp}} d + (1-\alpha_{\mathrm{hp}})x,\\
  x_{\mathrm{hp}} &= x - d,
\end{align}
where \(d\) is the internal drift estimate.

The result is then low-pass filtered to suppress high-frequency junk:
\begin{equation}
  b \leftarrow \alpha_{\mathrm{lp}} b + (1-\alpha_{\mathrm{lp}})x_{\mathrm{hp}}.
\end{equation}

The variable \(b\) is the prefiltered signal stored in the code as
\texttt{bp\_}. It is the signal actually used by the demodulator and loop.

\subsection{Power and RMS Estimate}

A smoothed power estimate is formed from the prefiltered signal:
\begin{equation}
  P \leftarrow \alpha_{\mathrm{pow}} P + (1-\alpha_{\mathrm{pow}})b^2,
\end{equation}
followed by
\begin{equation}
  \mathrm{RMS} = \sqrt{P + \varepsilon},
\end{equation}
with a small positive \(\varepsilon\) for numerical safety.

This RMS estimate is used for confidence, lock logic, and the optional coarse
acquisition stage.

% ---------------------------------------------------------------
\section{Synchronous Demodulation and Baseband States}

\subsection{Internal Oscillator}

The tracker maintains an internal oscillator phase \(\phi\). At each step,
the in-phase and quadrature carriers are
\begin{equation}
  c = \cos\phi,
  \qquad
  s = \sin\phi.
\end{equation}

\subsection{Demodulation}

The prefiltered signal \(b\) is synchronously demodulated as
\begin{equation}
  I_{\mathrm{meas}} = b\,\cos\phi,
  \qquad
  Q_{\mathrm{meas}} = -b\,\sin\phi.
\end{equation}

These are then low-pass filtered:
\begin{align}
  I &\leftarrow \alpha_{\mathrm{dem}} I + (1-\alpha_{\mathrm{dem}})I_{\mathrm{meas}},\\
  Q &\leftarrow \alpha_{\mathrm{dem}} Q + (1-\alpha_{\mathrm{dem}})Q_{\mathrm{meas}}.
\end{align}

The filtered pair \((I,Q)\) is the tracker’s baseband state.

\subsection{Amplitude and Coherence}

The code forms a tone-amplitude estimate from the demodulated baseband:
\begin{equation}
  A = 2\sqrt{I^2 + Q^2}.
\end{equation}

The factor of \(2\) compensates for the usual halving caused by synchronous
demodulation of a sinusoid.

A coherence-like score is then computed as
\begin{equation}
  \gamma
  =
  \operatorname{clamp}_{[0,1]}
  \left(
    \frac{A}{\sqrt{2}\,\mathrm{RMS}}
  \right).
\end{equation}

For a strongly dominant single tone with good preprocessing, this coherence
approaches \(1\). Broader-band or poorly matched signals push it downward.

% ---------------------------------------------------------------
\section{Confidence and Lock Logic}\label{sec:lock}

\subsection{Confidence Target}

The implementation combines an RMS score with the coherence score:
\begin{equation}
  r_{\mathrm{RMS}}
  =
  \operatorname{clamp}_{[0,1]}
  \left(
    \frac{\mathrm{RMS}}{\max(r_{\min},\varepsilon)}
  \right),
\end{equation}
\begin{equation}
  c^\star
  =
  \operatorname{clamp}_{[0,1]}(r_{\mathrm{RMS}}\,\gamma).
\end{equation}

This target confidence is then smoothed:
\begin{equation}
  c \leftarrow \alpha_{\mathrm{conf}} c + (1-\alpha_{\mathrm{conf}})c^\star.
\end{equation}

So confidence rises when the prefiltered signal has adequate RMS and the
demodulated baseband looks coherent.

\subsection{Hysteretic Lock Decision}

The tracker uses a hysteretic lock rule.

When currently unlocked, lock is declared only if
\begin{equation}
  \mathrm{RMS} \ge r_{\min}
  \quad\text{and}\quad
  \gamma \ge \gamma_{\mathrm{lock}}.
\end{equation}

When currently locked, unlock occurs if either
\begin{equation}
  \mathrm{RMS} < 0.7\,r_{\min}
\end{equation}
or
\begin{equation}
  \gamma < \gamma_{\mathrm{unlock}}.
\end{equation}

The two coherence thresholds satisfy
\begin{equation}
  0 \le \gamma_{\mathrm{unlock}} \le \gamma_{\mathrm{lock}} \le 1,
\end{equation}
so the lock decision has built-in hysteresis and does not chatter excessively.

% ---------------------------------------------------------------
\section{Phase Detector and PI Loop}\label{sec:loop}

\subsection{Phase Error}

The phase detector uses the low-pass baseband terms:
\begin{equation}
  e_\phi = \operatorname{atan2}(Q, I).
\end{equation}

To prevent excessive loop reactions during acquisition or disturbances, the
implementation clamps this phase error:
\begin{equation}
  e_\phi \leftarrow \operatorname{clamp}(e_\phi,-1.2,+1.2).
\end{equation}

The value returned by \texttt{getPhaseErrorRad()} is this clamped phase error.

\subsection{Loop Gains}

The code uses a type-II PLL-style structure with natural frequency
\begin{equation}
  \omega_n = 2\pi f_{\mathrm{loop}},
\end{equation}
proportional gain
\begin{equation}
  K_p = 2\zeta \omega_n,
\end{equation}
and integral gain
\begin{equation}
  K_i = \omega_n^2,
\end{equation}
where \(f_{\mathrm{loop}}\) is the configured loop bandwidth and \(\zeta\) is
the damping factor.

\subsection{Adaptation Gain}

The loop does not always use full authority. Instead, it scales the phase
error by a data-dependent adaptation gain:
\begin{equation}
  g =
  \begin{cases}
    c, & \text{if locked},\\[4pt]
    0.5\,c, & \text{if not locked}.
  \end{cases}
\end{equation}

So the loop reacts more cautiously when it is not yet confident.

\subsection{Frequency Integrator}

Before secondary corrections, the loop proposes
\begin{equation}
  \omega_{\mathrm{cand}}
  =
  \omega + K_i\,g\,e_\phi\,\Delta t.
\end{equation}

This is the integral part of the PLL/FLL-style update.

\subsection{Recentering When Confidence Is Very Poor}

If the confidence is very low,
\begin{equation}
  c < 0.10,
\end{equation}
the code slowly pulls the frequency back toward the configured nominal
frequency \(\omega_{\mathrm{nom}}\):
\begin{equation}
  \omega_{\mathrm{cand}}
  \leftarrow
  \omega_{\mathrm{cand}}
  +
  \beta_{\mathrm{rec}}
  (\omega_{\mathrm{nom}} - \omega_{\mathrm{cand}}),
\end{equation}
where
\begin{equation}
  \beta_{\mathrm{rec}} = 1 - \exp\!\left(-\frac{\Delta t}{\tau_{\mathrm{rec}}}\right).
\end{equation}

This is a practical recovery feature, not part of a textbook PLL.

\subsection{Slew Limit and Frequency Band Limit}

The implementation limits how quickly the frequency may move. If the maximum
slew rate is \(f'_{\max}\) in \si{Hz/s}, then
\begin{equation}
  |\Delta\omega| \le 2\pi f'_{\max}\Delta t.
\end{equation}

After this rate limit, the frequency is clamped to the configured search band:
\begin{equation}
  2\pi f_{\min} \le \omega \le 2\pi f_{\max}.
\end{equation}

\subsection{Phase Advance}

After the updated oscillator frequency is obtained, the internal phase is
advanced as
\begin{equation}
  \phi \leftarrow \phi + (\omega + K_p g e_\phi)\Delta t,
\end{equation}
and then wrapped into \([-\pi,\pi]\).

Thus the returned phase is the internal oscillator phase, not a separately
estimated analytic signal phase.

% ---------------------------------------------------------------
\section{Output Frequency Smoothing}

The raw tracked frequency is simply
\begin{equation}
  f_{\mathrm{raw}} = \frac{\omega}{2\pi}.
\end{equation}

The public frequency output is a smoothed version:
\begin{equation}
  f_{\mathrm{smooth}}
  \leftarrow
  \alpha_{\mathrm{out}} f_{\mathrm{smooth}}
  +
  (1-\alpha_{\mathrm{out}}) f_{\mathrm{raw}}.
\end{equation}

The returned value from \texttt{getFrequencyHz()} is \(f_{\mathrm{smooth}}\),
while \texttt{getRawFrequencyHz()} returns \(f_{\mathrm{raw}}\).

This separation is helpful: the raw loop state reacts more quickly, while the
public output is less jittery.

% ---------------------------------------------------------------
\section{Optional Coarse Acquisition Assist}\label{sec:coarse}

The implementation optionally includes a simple threshold-based coarse
acquisition aid.

\subsection{Thresholding Against Current RMS}

After the internal substeps are complete, the tracker uses the final
prefiltered sample \(b\) and the current RMS to form a hysteretic threshold
\begin{equation}
  T_c = \max(\eta\,\mathrm{RMS}, 10^{-9}),
\end{equation}
where \(\eta\) is the configured coarse hysteresis fraction.

The coarse detector first arms when
\begin{equation}
  b \le -T_c,
\end{equation}
and later records an event when the signal rises to
\begin{equation}
  b \ge +T_c.
\end{equation}

So it looks for a sufficiently negative excursion followed by a sufficiently
positive excursion.

\subsection{Period from Successive Events}

Let the times of successive coarse events be \(t_k\). Then the measured period
is
\begin{equation}
  T_k = t_k - t_{k-1},
\end{equation}
and the measured coarse frequency is
\begin{equation}
  f_{c,k} = \frac{1}{T_k}.
\end{equation}

Only frequencies inside the configured band are accepted.

Accepted measurements are exponentially smoothed:
\begin{equation}
  f_c \leftarrow \alpha_c f_c + (1-\alpha_c) f_{c,k}.
\end{equation}

If no fresh coarse event arrives for too long, the coarse estimate times out
and is marked invalid.

\subsection{Gentle Pull Toward the Coarse Estimate}

When coarse assist is enabled and valid, the loop gently pulls the candidate
frequency toward the coarse estimate:
\begin{equation}
  \omega_{\mathrm{cand}}
  \leftarrow
  \omega_{\mathrm{cand}}
  +
  g_c \beta_c(\omega_c - \omega_{\mathrm{cand}}),
\end{equation}
where \(\omega_c = 2\pi f_c\), and
\begin{equation}
  \beta_c = 1 - \exp\!\left(-\frac{\Delta t}{\tau_{\mathrm{pull}}}\right).
\end{equation}

The assist gain is stronger when confidence is poor:
\begin{equation}
  g_c =
  \begin{cases}
    0.15(1-c), & \text{if locked},\\[4pt]
    1.0(1-c), & \text{if not locked}.
  \end{cases}
\end{equation}

This means coarse assist mainly helps acquisition and recovery, while staying
gentle once the loop is already well locked.

% ---------------------------------------------------------------
\section{Numerical Safeguards and Implementation Notes}

Several practical safeguards are built into the code:

\begin{itemize}[leftmargin=1.4em]
  \item All one-pole coefficients are formed with bounded exponential
        functions.
  \item Very large public \(\Delta t\) values are split into smaller stable
        internal steps.
  \item Frequency slew rate is explicitly limited.
  \item Frequency is always clamped to the physical search band.
  \item The phase detector output is bounded.
  \item Stable \(\operatorname{hypot}\)-style amplitude calculation is used for
        \(\sqrt{I^2+Q^2}\).
  \item Configuration values are sanitized before use.
\end{itemize}

These choices make the implementation more robust than a minimal textbook PLL.

\begin{table}[t]
  \centering
  \caption{Main internal states}
  \label{tab:states}
  \small
  \begin{tabular}{@{}lll@{}}
    \toprule
    Quantity & Code name & Meaning \\
    \midrule
    \(\omega\) & \texttt{omega\_rad\_s\_} & internal oscillator frequency \\
    \(\phi\) & \texttt{phase\_rad\_} & internal oscillator phase \\
    \(e_\phi\) & \texttt{phase\_error\_rad\_} & clamped phase error \\
    \(d\) & \texttt{dc\_lp\_} & drift estimate for HP stage \\
    \(b\) & \texttt{bp\_} & prefiltered signal \\
    \(I,Q\) & \texttt{i\_lp\_}, \texttt{q\_lp\_} & baseband states \\
    \(\mathrm{RMS}\) & \texttt{rms\_} & signal RMS estimate \\
    \(A\) & \texttt{amplitude\_} & tone-amplitude estimate \\
    \(\gamma\) & \texttt{coherence\_} & coherence score \\
    \(c\) & \texttt{confidence\_} & smoothed confidence \\
    \bottomrule
  \end{tabular}
\end{table}

% ---------------------------------------------------------------
\section{Summary}

The implemented \texttt{PLLFreqTracker} is a practical one-frequency tracker
with the following structure:
\begin{itemize}[leftmargin=1.4em]
  \item front-end drift removal and broad bandpass preprocessing,
  \item synchronous demodulation into low-pass baseband \(I/Q\),
  \item phase error from \(\operatorname{atan2}(Q,I)\),
  \item bounded PI loop on oscillator frequency,
  \item confidence and hysteretic lock logic,
  \item optional coarse acquisition assistance,
  \item explicit slew limits, band limits, and variable-\(dt\) handling.
\end{itemize}

It is compact, numerically cautious, and well suited to embedded tracking of a
dominant narrowband wave component.

\end{document}
